\clearpage
\fancyhead[L]{\small JUNO.}
\hypertarget{native-265-lower}{}
\begin{center}
JUNO.

\rule{0.6\linewidth}{0.4pt}

Monatliche Correspondenz. Band XI. S. 184–193. Februar 1805.

\rule{0.6\linewidth}{0.4pt}
\end{center}

—\quad—\quad—

Wie immer lieferte auch diesmal Dr. \textsc{Gauss}, der mit ausdauerndem Fleiß weiterarbeitet, die interessantesten und wichtigsten Beiträge. Deshalb beginnen wir auch hier unseren Bericht über die neuesten Beobachtungen und Berechnungen zur Juno mit seinen Beiträgen. Im November 1804 hatte dieser Astronom vergeblich darauf gewartet, dass man ihm gute Meridianbeobachtungen der Juno mitteilt, mit deren Hilfe er ihre Elemente weiter verbessern könnte. Da er jedoch bis zum Dezember keine erhielt und die nach den dritten Elementen berechneten Positionen um einige Minuten von den beobachteten abwichen, stützte er die folgenden IV. Elemente auf mehrere Beobachtungen der Juno, die er teils selbst am Kreismikrometer vorgenommen, teils von Dr. \textsc{Olbers} erhalten hatte.

\begin{center}
\emph{IV. Elemente der Junobahn.}
\end{center}

\begin{center}
\begingroup\fontsize{9}{11}\selectfont\setlength{\tabcolsep}{3pt}\renewcommand{\arraystretch}{1.1}
\begin{tabular}{lr}
Epoche 1805 . . . . . . . . . & $42^\circ\ 41'\ 34''$ \\
Aphel . . . . . . . . . & $233\ 23\ 47$ \\
Aufsteigender Knoten . . . . . . & $171\ 4\ 12$ \\
Neigung . . . . . . . . . . . & $13\ 4\ 9$ \\
Tägliche Bewegung . . . . . . . . & $812''091$ \\
Logarithmus der Halbachse . . . . & $0.426935$ \\
Exzentrizität . . . . . . . . . . & $0.256841$ \\
\end{tabular}\endgroup
\end{center}

\indent Aus diesen Elementen berechnete Dr. \textsc{Gauss} ihren geozentrischen Lauf für Dezember 1804 sowie Januar und Februar 1805. Die Ephemeride für den mittleren Monat haben wir schon im vorigen Heft mitgeteilt. Daher folgen hier nur die Ephemeriden für Dezember und Februar [Diese Ephemeriden sind in diesem Werk nicht abgedruckt.] —\quad—\quad—

Dr. \textsc{Gauss} machte im Dezember und Januar die folgenden drei Beobachtungen.

\begin{center}
\begingroup\fontsize{9}{11}\selectfont\setlength{\tabcolsep}{3pt}\renewcommand{\arraystretch}{1.1}
\begin{tabular}{l|r|r|r}
& Mittlere Zeit & Rekt. $\JunoSymbol$ & \shortstack{Declinatio\\australis}\\ \hline
1804, 29. Dezember & $6^{\mathrm u}\ 53^{\mathrm m}\ 43^{\mathrm s}$ & $7^\circ\ 23'\ 11''$ & $7^\circ\ 28'\ 2''$ \\
30 & $6\ 19\ 50$ & $7\ 44\ 5$ & $7\ 19\ 57$ \\
1805, 5. Januar & $5\ 6\ 57.5$ & $9\ 59\ 42$ & $6\ 30\ 56$ \\
\end{tabular}\endgroup
\end{center}
\par\Needspace{11\baselineskip}
\SourceLeaf{266}\par
\indent Er schrieb uns dazu, dass ihm die letzte Beobachtung für sich genommen zwar gut erscheine. Der Vergleich mit einem Stern aus der Histoire céleste mache sie jedoch sehr zweifelhaft. Am 5. Januar verglich Dr. \textsc{Gauss} Juno achtmal mit dem Stern achter Größe. Seine Position ist in der Histoire céleste auf S. 135 wie folgt angegeben.

\begin{center}
\begingroup\fontsize{10}{12}\selectfont\setlength{\tabcolsep}{3pt}\renewcommand{\arraystretch}{1.1}
\begin{tabular}{c|r|r}
& Rekt. & Zenitdistanz\\ \hline
8 & $0^{\mathrm u}\ 40^{\mathrm m}\ 14^{\mathrm s}$ & $55^\circ\ 19'\ 15''$ \\
\end{tabular}\endgroup
\end{center}


\noindent Doch schon nach bloßem Augenmaß schien ihm dieser Stern mit drei anderen, die über ihm standen, eine andere rhomboidale Figur zu bilden, als sich aus den Angaben in der Histoire céleste ergab. Da Dr. \textsc{Gauss} die Position dieses Sterns genauer bestimmt haben wollte, nutzten wir dazu den ersten klaren Abend. Seine Vermutung bestätigte sich vollständig. Wir erhielten als südliche Deklination nicht $6^\circ\,28'$, sondern $6^\circ\,23'\,9''$. Danach würde die am 5. Januar beobachtete Deklination der Juno $6^\circ\,25'\,47''$ betragen. Aus den IV. Elementen ergeben sich diese Werte.

\begin{center}
\begingroup\fontsize{10}{12}\selectfont\setlength{\tabcolsep}{3pt}\renewcommand{\arraystretch}{1.1}
\begin{tabular}{l|r|r}
& Rekt. $\JunoSymbol$ & Deklination\\ \hline
1804, 30. Dezember & $7^\circ\ 44'\ 14''$ & $7^\circ\ 20'\ 16''$ \\
1805, 5. Januar & $9\ 59\ 44$ & $6\ 25\ 58$ \\
\end{tabular}\endgroup
\end{center}


\noindent Nach der letzten Reduktion stimmen die beobachteten Positionen mit diesen Werten sehr gut überein. Zu unserer Bestimmung der Position dieses zweifelhaften Sterns müssen wir jedoch anmerken, dass es bei seiner Kulmination noch so hell war, dass diese Beobachtungen kaum als völlig zuverlässig gelten können. —\quad—\quad—

[Mit diesen IV. Elementen berechnete \textsc{Gauss} den Zodiakus der Juno nach den oben auf \hyperlink{native-111}{S. 106} abgedruckten Vorschriften, Mon. Corr. Bd. XI, S. 225–228. März 1805.]

\hfill [\textsc{von Zach.}]

